\documentclass{article}
\usepackage[utf8]{inputenc}
\usepackage[T1]{fontenc}
\usepackage{geometry}
\geometry{
    a4paper,
    left=2.5cm,
    right=2.5cm,
    top=2.5cm,
    bottom=2.5cm
}
\usepackage{graphicx}
\usepackage{float}
\usepackage{hyperref}
\usepackage[english]{babel}
\usepackage{amsmath}
\usepackage{amssymb}
\usepackage{array}
\usepackage{booktabs}
\usepackage{caption}
\usepackage{subcaption}
\usepackage{enumitem}
\usepackage{fancyhdr}
\usepackage{titlesec}
\usepackage{url}
\usepackage{textcomp}
\usepackage{xcolor}
\usepackage{listings}
\usepackage{multirow}
\usepackage{longtable}
\usepackage{tabularx}
\usepackage{times}

% Set up page style
\pagestyle{fancy}
\fancyhf{}
\rfoot{\thepage}
\renewcommand{\headrulewidth}{0pt}

% Title formatting
\titleformat{\chapter}[display]{\normalfont\bfseries}{\chaptertitle\thechapter}{0pt}{\Huge}
\titlespacing*{\chapter}{0pt}{50pt}{40pt}

% Document info
\title{Template Document}
\author{Academic Support Agent}
\date{\today}

\begin{document}
\maketitle

\pagenumbering{gobble} % keep title page without a number
\maketitle



\pagenumbering{arabic} % start page numbers again
\setcounter{page}{1}

\section{Examples of LaTeX coding for future use}
\subsection{To keep hierarchy}
\subsubsection{Paragraph writing}

\paragraph{This is a paragraph}
\lipsum[2]

\subparagraph{And this is a subparagraph}
\lipsum[3]

\subsection{\textit{Not indented}}

\noindent Starts as left as possible Not indented Not indented Not indented

\section{Highlighting and footnotes}
You can make words \textbf{bold}, \textit{italicise} them, \underline{underline words} or \textbf{make them \emph{stand out} regardless of the surrounding}. You can break a line\\ mid sentence and make footnotes like this \footnote{a footnote}.

\section{Equations}

% Forgot a symbol? See the LaTeX symbol wiki https://en.wikipedia.org/wiki/Wikipedia:LaTeX_symbols

\subsection{As part of text}

In total 85 distinct galaxies were identified in the Hubble Deep Field image provided, the list of which can be found in Appendix B at the end of this report. Poisson statistics states that the error in the number of galaxies counted ($N_x$) is simply the root of the count ($\sqrt{N_x}$), this can be represented as a percentage by the following equation, $N_x^2 + N^2$.

\subsection{In the middle, not numbered} 

$$ \alpha \beta \times \frac{2G}{x^2B_n}$$

$$ \frac{-b \pm \sqrt{b^2-4ac}}{2a} + \sin{x}$$

% mag = 5 * np.log10(dlpc/10) + M
$$ \mathrm{mag} = 5 \times \mathrm{log}_{10}\left(\frac{D_{l}}{10}\right) + \mathrm{M_{corr}}$$

\subsubsection{And in the middle, numbered}
\begin{equation}
    A=bx+24\times F_x
    \label{eq:lovely formula}
\end{equation}

\section{Adding images}
Hello, here is some text without a meaning.  This text should show what a printed text will look like at this place. If you read this text, you will get no information.  Really?  Is there no information?  Is there.

\subsection{Lonely image}
\begin{figure}[h!]
    \centering % keeps the image in the centre
    \includegraphics[width=0.5\textwidth]{snailfrog.png}
    \caption{A visualisation of a frog.}
    \label{fig:frog}
\end{figure}

\lipsum[7]

\subsubsection{Images side by side}
\begin{figure}[h]
    \centering
    \begin{minipage}{0.45\textwidth}
        \centering
        \includegraphics[width=0.9\textwidth]{frog.jpg}
        \caption{first figure}
    \end{minipage}\hfill
    \begin{minipage}{0.45\textwidth}
        \centering
        \includegraphics[width=0.9\textwidth]{serious_frog.jpeg}
        \caption{second figure}
    \end{minipage}
\end{figure}

\section{Citing and referencing}

\subsection{Referencing figures and equations}
 Expression $4\times 3=G\times x$ naturally follows from Eq \ref{eq:lovely formula}, and both of these things have a lot to do with Fig \ref{fig:frog in multicol}.
 
 \subsection{Citing a paper}
This statement has a citation at the end of it \cite{toadetal1958}, and this one has two \cite{toadetal1958, squeaker1982}. A citation with parenthesis is sure to follow \citep{siddiqi2004interspecific}.

\section{Tables}

\begin{table}[ht]
\centering                      % used for centering table
\begin{tabular}{c c c c}        % centered columns (4 columns)
Case & Ar & H & He \\ [0.5ex]   % inserts table %heading
\hline                          % inserts 1 horizontal line
1 & 50 & 837 & 970 \\           % table contents
2 & 47 & 877 & 230 \\                         
3 & 31 & 25  & 415 \\
4 & 35 & 144 & 2356 \\
5 & 45 & 300 & 556 \\ [1ex]     % [1ex] adds vertical space
\end{tabular}
\caption{Imaginary experiment results.} 
\label{table:nonlin}            % is used to refer this table in the text
\end{table}
 
\section{Listing things}
 
\subsection{Un-ordered list}
 \begin{itemize}
	\item One
	\item Two
	\item Three
\end{itemize}

\subsection{Ordered and nested lists}
\begin{enumerate}
	\item One
    \begin{enumerate}
    	\item Two
        \item Three
        \item Four
    \end{enumerate}
    \item Five
    \item Six
\end{enumerate}

\section{Example of symbol in section heading below}
\subsection{Lyman-\texorpdfstring{$\alpha$}{Lg}}

\section*{Example of an unnumbered section}
It won't appear in the table of contents

\section{Columns}
\begin{multicols}{2}
\subsection{Two Columns}
\lipsum[1-2]

\subsubsection{A picture in a multi-column environment}

\begin{figure}[H]
    \includegraphics[width=\linewidth]{serious_frog.jpeg}
    \caption{A serious frog.}
    \label{fig:frog in multicol}
\end{figure}

\end{multicols}


\begin{multicols}{3}[\subsection{Three Columns}]
\lipsum[4-6]
\end{multicols}

\newpage


%~~~~~~~~~ Bibliography

\setcitestyle{authoryear}
\end{document}